\begin{itemize}
    \item ¿Por qué no basta con "tener Wi-Fi"?
    Una solución IoT completa involucra múltiples capas de comunicación. Algunas se encargan de conectar sensores al microcontrolador, otras permiten enlazar el nodo con internet, y otras definen cómo se intercambian los datos a nivel de aplicación. Reducir el IoT únicamente a la conectividad Wi-Fi constituye un error conceptual frecuente, ya que ignora la complejidad real de los sistemas embebidos conectados.

    \begin{itemize}
        \item Comunicación interna del nodo embebido (dentro del hardware)
        Antes de transmitir información a la red, el sistema debe capturar y procesar correctamente los datos provenientes de los sensores:

        \begin{itemize}
            \item GPIO y señales digitales: lectura de estados y activación de actuadores.
            \item Entradas analógicas y ADC: adquisición de variables continuas como temperatura o presión.
            \item PWM: control de velocidad, intensidad luminosa o posición.
            \item I2C: comunicación con múltiples sensores usando pocos conductores.
            \item SPI: transmisión de alta velocidad con memorias, pantallas o sensores especializados.
            \item UART: comunicación serial para depuración o conexión con módulos externos.
        \end{itemize}

        \item Tecnologías inalámbricas para conectar el nodo a la red

        La selección de la tecnología de comunicación depende del equilibrio entre alcance, consumo energético, velocidad de transmisión y costo de implementación. En la Tabla~\ref{tab:Tabla_4_PDCIoT} se resumen las tecnologías más utilizadas.

        \begin{table*}
    

\centering
\renewcommand{\arraystretch}{1.3}
\label{tab:Tabla_4_PDCIoT}
\begin{tabular}{|l|c|c|c|p{2cm}|}
\hline
\textbf{Tecnología} & \textbf{Alcance} & \textbf{Consumo} & \textbf{Tasa de datos} & \textbf{Uso típico} \\
\hline
Wi-Fi & $\sim$100 m & Alto & Hasta 1 Gbps & Hogares, oficinas, prototipos ESP32 \\
\hline
Bluetooth / BLE & 10--100 m & Bajo & 1--3 Mbps & Wearables, sensores cercanos \\
\hline
ZigBee & $\sim$100 m & Bajo & 250 kbps & Domótica, automatización industrial \\
\hline
LoRa / LPWAN & 2--15 km & Muy bajo & 0.3--50 kbps & Agricultura, ciudades inteligentes, campo abierto \\
\hline
NFC / RFID & $<$ 10 cm & Muy bajo & 424 kbps & Identificación, pagos, logística \\
\hline
Celular 4G/5G & Km & Variable & Hasta Gbps & Vehículos, monitoreo remoto sin red local \\
\hline
Ethernet & Cableado & Bajo/medio & Alta & Ambientes industriales fijos, alta estabilidad \\
\hline
\end{tabular}
\caption{Tecnologías inalámbricas para conectar el nodo a la red}
\label{tab:tecnologias_inalambricas}
\end{table*}[h]

        No existe una tecnología universalmente superior; la elección depende del caso de uso específico.

        \item Notas específicas de cada tecnología

        \begin{itemize}
            \item Wi-Fi: basado en el estándar 802.11, ideal para transmisión de grandes volúmenes de datos y conexión directa a internet.
            \item Bluetooth / BLE: optimizado para bajo consumo, fundamental en dispositivos portátiles y sensores de batería.
            \item ZigBee: diseñado para redes de sensores distribuidos que operan de forma cooperativa.
            \item Thread: protocolo mallado basado en IPv6, auto-reparable y eficiente energéticamente.
            \item LoRaWAN: orientado a telemetría de largo alcance con bajo consumo y baja frecuencia de transmisión.
            \item Celular / NB-IoT: permite independencia de redes locales, aunque con mayor complejidad operativa.
        \end{itemize}

        \item Protocolos de aplicación: cómo se empaqueta y transporta la información

        Una vez establecida la conexión de red, es necesario definir el protocolo de intercambio de datos:

        \begin{itemize}
            \item HTTP / REST: modelo cliente-servidor basado en solicitudes y respuestas. Adecuado para servicios web y pruebas iniciales.
            \item MQTT: protocolo ligero basado en publicador/suscriptor, altamente eficiente para telemetría y sistemas distribuidos. El broker gestiona la distribución de mensajes.
            \item WebSocket: permite comunicación bidireccional continua mediante conexión persistente, útil para dashboards y control en tiempo real.
        \end{itemize}

        En aplicaciones con ESP32, MQTT resulta especialmente formativo para comprender la lógica IoT, mientras que HTTP facilita la integración con servicios web.

        \item Criterio de selección resumido

        No existe una solución única óptima. La elección adecuada surge del análisis de los requerimientos del sistema:

        \begin{itemize}
            \item ¿Qué alcance de comunicación se necesita?
            \item ¿Qué volumen de datos se transmitirá y con qué frecuencia?
            \item ¿Cuál es el consumo energético permitido?
            \item ¿Cuál es el costo y la facilidad de implementación en el entorno específico?
        \end{itemize}

    \end{itemize}
\end{itemize}